\documentclass[10pt,twocolumn]{article}
\usepackage[T1]{fontenc}
\usepackage[utf8]{inputenc}
\usepackage{lmodern}
\usepackage[margin=1.8cm,top=2.4cm,bottom=2cm,columnsep=0.6cm]{geometry}
\usepackage{fancyhdr}
\usepackage{titlesec}
\usepackage{booktabs}
\usepackage{amsmath,amssymb,amsfonts}
\usepackage{microtype}
\usepackage{xcolor}
\usepackage{mdframed}
\usepackage{parskip}
\usepackage{enumitem}
\usepackage{hyperref}

\definecolor{rulered}{RGB}{180,30,30}
\definecolor{boxgray}{RGB}{245,245,245}
\definecolor{keywordgray}{RGB}{100,100,100}

\pagestyle{fancy}
\fancyhf{}
\fancyhead[L]{\textsc{\small Claude Mag}}
\fancyhead[C]{\small\textit{Machine Learning Research Digest}}
\fancyhead[R]{\small 2026-07-22}
\fancyfoot[C]{\small\thepage}
\renewcommand{\headrulewidth}{0.8pt}

\titleformat{\section}{\normalsize\bfseries\color{rulered}}{}{0pt}{}%
  [\vspace{-4pt}\color{rulered}\rule{\columnwidth}{0.4pt}]
\titlespacing{\section}{0pt}{8pt}{4pt}

\hypersetup{colorlinks=true,urlcolor=rulered,linkcolor=black}

\begin{document}
\setlength{\parindent}{0pt}
\setlength{\parskip}{4pt}

{\small\color{keywordgray}\textbf{Keywords:}
  reinforcement learning \textbullet{}
  compositional reasoning \textbullet{}
  post-training \textbullet{}
  skill composition \textbullet{}
  rewrite grammars}\par\vspace{4pt}

{\LARGE\bfseries\raggedright
  RL Post-Training Builds Compositional Reasoning Strategies}\par\vspace{4pt}

{\large\itshape\raggedright
  UCL Gatsby Researchers Demonstrate That RL with Only Binary Reward Can
  Compose Pretrained Primitive Skills Into New Higher-Level Strategies
  That Rejection Fine-Tuning Cannot Reach}\par\vspace{6pt}

{\small\textbf{By} Azwar Abdulsalam, Nishil Patel, Andrew Saxe
  (UCL Gatsby Computational Neuroscience Unit)}\par\vspace{2pt}

{\small\color{keywordgray}arXiv:2607.07646 \textbullet{}
  July 8, 2026}
\par\vspace{2pt}\noindent{\color{rulered}\rule{\columnwidth}{0.6pt}}\vspace{6pt}

A central question in post-training research is whether reinforcement
learning merely amplifies skills already latent in a pretrained base
model, or whether it can build genuinely new capabilities by composing
those primitives in ways the model has never seen.
Researchers at UCL's Gatsby Computational Neuroscience Unit now
answer this question decisively using a controlled symbol-rewrite
grammar environment: RL with binary final-answer reward not only
outperforms rejection fine-tuning on composite tasks, it produces
a new compositional mechanism that emerges in two identifiable phases.

\begin{mdframed}[backgroundcolor=boxgray,linecolor=rulered,linewidth=1pt,
    innerleftmargin=6pt,innerrightmargin=6pt,
    innertopmargin=6pt,innerbottommargin=6pt]
\textbf{\small AT A GLANCE}\par\vspace{4pt}
\begin{description}[leftmargin=0pt,labelsep=4pt,itemsep=2pt,parsep=0pt,
    font=\small\bfseries]
  \item[Problem:] It is unknown whether RL post-training amplifies
    existing skills or composes them into genuinely new strategies
    unavailable to the base model.
  \item[Why it matters:] If RL can only amplify, it is bounded by
    pretraining; if it can compose, it can transcend pretraining
    data in principled ways.
  \item[Method:] Controlled rewrite-grammar environment where
    primitives and compositions are known; Transformer pretrained
    on primitives, RL post-trained on composite tasks with binary reward.
  \item[Result:] RL reaches 78\% on composite problems vs.\ 40\% for
    RFT plateau; compositional strategies emerge in two phases.
  \item[Evidence:] Phase transition at step $\sim$600; sequential
    compositions emerge first, then parallel compositions.
\end{description}
\end{mdframed}

\section{The Big Picture}

The dominant post-training recipe for reasoning models involves
either rejection fine-tuning (RFT) --- filtering model rollouts for
correct answers and training on them --- or reinforcement learning
(RL) with verifiable reward.
These two approaches are widely believed to have different capability
profiles: RFT is seen as a data-efficient curriculum tool that
stays within the distribution of the base model, while RL is
believed to discover new strategies through trial and error.
But the evidence for this distinction has been largely anecdotal.

The Gatsby team designs a setting where ground truth is available
for exactly which strategies the model uses at each step, making it
possible to verify mechanistically whether RL produces new
compositional behavior or merely amplifies existing primitives.

\section{Why Existing Approaches Fall Short}

Both existing post-training approaches have limitations that are
difficult to diagnose on open-ended language tasks.
RFT cannot improve beyond the base model's rollout distribution:
if the base model rarely produces composite solutions, RFT has
no signal to amplify.
RL can in principle explore beyond the base distribution, but
it is unclear whether the gradient signal from binary reward is
sufficient to wire together primitive capabilities into structured
compositions, or whether it simply boosts the probability of
already-composite rollouts that happen to appear at low frequency.

In language settings, these mechanisms are entangled with content
and fluency; the rewrite-grammar setting surgically isolates
the compositional question from all other confounds.

\section{Core Mechanism}

EcoSpec's RL training is shown to operate through a two-phase
compositional mechanism.
In \textbf{Phase 1}, RL sharpens the model's execution of individual
primitive rules: the probability of applying the correct primitive
to the correct symbol increases, reducing errors in primitive steps.
In \textbf{Phase 2}, the model begins chaining these primitive
applications --- first sequentially (applying rule $r_i$ then $r_j$
in a coordinated trace), then in parallel (applying $r_i$ and $r_j$
to independent subterms simultaneously within a single generation).
The transition between phases is sharp and reproducible across seeds.

\section{Technical Formulation}

The rewrite grammar defines a set of primitive rules
$\mathcal{R} = \{r_1, \ldots, r_K\}$.
Each primitive $r_i$ rewrites a specific input symbol pattern
$s_{\text{in}}^{(i)}$ to $s_{\text{out}}^{(i)}$.
A composite problem of arity 2 requires applying rules
$r_i$ and $r_j$ either sequentially or in parallel:
\[
\text{Sequential:}\quad r_j(r_i(x)) = y
\quad
\text{Parallel:}\quad (r_i \!\parallel\! r_j)(x_1, x_2) = (y_1, y_2)
\]
The Transformer is pretrained only on single-rule traces (input, output
pairs with $r_i$ applied once). At post-training, it receives composite
problem inputs and binary reward for producing the correct final output.

The policy gradient update under REINFORCE takes the form:
\[
\nabla_\theta \mathcal{J}(\theta)
= \mathbb{E}_{\tau \sim \pi_\theta}
  \left[ R(\tau) \sum_{t=1}^{T} \nabla_\theta \log \pi_\theta(a_t \mid s_t) \right]
\]
where $R(\tau) \in \{0,1\}$ is the binary correctness reward for
the full trace $\tau$ and the sum is over all generation steps.

\section{Learning Procedure}

\begin{enumerate}[itemsep=2pt,parsep=0pt]
  \item \textbf{Pretraining:} The Transformer is trained on a corpus
    of single-rule rewrite traces covering all $K$ primitives uniformly.
    After pretraining, the model reliably applies any single primitive
    but has near-zero accuracy on composite tasks.
  \item \textbf{RL post-training:} REINFORCE with binary reward is
    applied to composite tasks.
    No intermediate-step reward, no step annotations.
    Training runs for 2000 gradient steps across all conditions.
  \item \textbf{RFT baseline:} Rollouts from the pretrained model on
    composite tasks are filtered for correctness; the model is fine-tuned
    on the filtered set.
    RFT training uses the same FLOPs as RL where possible.
\end{enumerate}

\section{Experimental Findings}

\begin{center}
\small
\begin{tabular}{lcc}
\toprule
Method & Accuracy & Plateau \\
\midrule
Base model (pass@1) & 4.3\% & --- \\
Base model (pass@100) & 19.1\% & --- \\
RFT & 40.2\% & Step 400 \\
RL (this work) & \textbf{78.3\%} & Step 2000 \\
\bottomrule
\end{tabular}
\end{center}

The phase transition is visible in trace analysis: the rate of
sequential compositions rises from near-zero to $\sim$30\% by
step 600, then stabilizes while parallel compositions rise from
near-zero to $\sim$45\% by step 2000.

The base model's pass@100 rate (19.1\%) establishes a ceiling for
RFT: RFT cannot learn from solutions the model never produces under
any sampling strategy, yet RFT still substantially surpasses this
baseline --- suggesting that rare composite rollouts provide
sufficient training signal when filtered correctly.

\section{Ablations and Interpretation}

\textbf{Sequential before parallel:} Parallel compositions never
emerge before sequential ones across any random seed or grammar
instantiation.
The authors interpret this as a task-induced curriculum:
sequential composition is a prerequisite for parallel composition
because it requires the model to coordinate rule application across
a linear chain before coordinating across independent subterms.

\textbf{Binary vs.\ step-level reward:} Providing step-level
intermediate reward speeds up phase 1 (primitive sharpening)
but does not accelerate phase 2 (compositional emergence),
suggesting that compositional emergence is driven by global
reward signal optimizing the full trace rather than local
per-step feedback.

\textbf{RFT saturation:} RFT saturates because the base model's
composite-solution rollout rate is low, providing a thin training
signal.
Once that signal is exhausted, gradient updates from the sparse
correct rollouts no longer produce the global policy restructuring
that RL achieves through exploration.

\textbf{Implications:} These results support a view of RL as a
\emph{composer} rather than merely an \emph{amplifier}: given
primitives in pretraining and composition problems at post-training,
RL can discover how to chain those primitives into strategies
that the base model has never executed.

\section*{Reference}

Azwar Abdulsalam, Nishil Patel, Andrew Saxe.\
\textbf{RL Post-Training Builds Compositional Reasoning Strategies}.
arXiv:2607.07646, July 2026.
\url{https://arxiv.org/abs/2607.07646}

\end{document}
